\documentclass[12pt,oneside]{article}
\usepackage[margin=2cm]{geometry}
\usepackage{multirow}
\usepackage{graphicx}
\usepackage[colorlinks=true,citecolor=black,urlcolor=blue]{hyperref}

%\newcommand{\embed}[5]{\begin{center}\fbox{\includegraphics[trim=#1 #2 #3 #4,clip,width=\textwidth]{#5}}\end{center}\clearpage}
\newcommand{\embed}[5]{\begin{center}\includegraphics[trim=#1 #2 #3 #4,clip,width=\textwidth]{#5}\end{center}\clearpage}

\usepackage{fancyhdr}
\setlength\headheight{2cm}
\fancyhead{}
\fancyfoot{}
\renewcommand{\headrulewidth}{1pt}
\renewcommand{\footrulewidth}{0pt}
\lhead{Bioinformatics Open Source Conference (BOSC) 2014\\Boston}
\rhead{\includegraphics[height=1cm]{../../logos/100px-Pear.png}}

\newcommand{\session}[3]{\clearpage\lhead{BOSC 2014 -- \textit{#1}\\#2\\#3}\newpage}

\title{%Hack to get the logo on the PDF front page:
\vspace{4cm}
\includegraphics[width=4cm]{../../logos/100px-Pear.png}\\
%Hack to get some white space using a blank line:
~\\Bioinformatics Open Source Conference (BOSC)\\
Talk Abstracts (Draft)}

\date{11 \& 12 July 2014, Boston, MA}

\begin{document}
\pagestyle{fancy}

\maketitle

\noindent For the latest schedule, please see:\\
\url{http://www.open-bio.org/wiki/BOSC_2014_Schedule}

TODO - Sponsors and logos!

\thispagestyle{empty}

\session{CodeFest}{Pre-BOSC, 9 \& 10 July 2014}{Organisers: Brad Chapman \& Michael Heuer}

\noindent
As in recent years, for the two days before BOSC we are holding the
BOSC CodeFest, an informal ``Coding Festival'' or mini-hackathon:
\url{http://www.open-bio.org/wiki/Codefest_2014}

\begin{center}
\includegraphics[width=4cm]{../../logos/Hack-reduce-logo.png}
\end{center}

\noindent
The BOSC CodeFest 2014 is being hosted by \href{http://www.hackreduce.org/}{hack/reduce},
a wonderful hacker space in Cambridge.

\begin{center}
\includegraphics[width=5cm]{../../logos/150px-Aws-logo.jpeg}
\end{center}

\noindent
Thanks to \href{http://aws.amazon.com/}{Amazon Web Services} all participants will receive a $\$100$
AWS credit to support work at the hackathon.

\begin{center}
\begin{tabular}{cc}
\multirow{2}{*}{\raisebox{\dimexpr-\height+\ht\strutbox}{\includegraphics[height=3cm]{../../logos/Harbinger-Partners.png}}} &
\raisebox{\dimexpr-\height+\ht\strutbox}{\includegraphics[width=4cm]{../../logos/Curoverse.png}} \\
& \raisebox{\dimexpr-\height+\ht\strutbox}{\includegraphics[width=4cm]{../../logos/Arvados.png}}
\end{tabular}
\end{center}

\noindent
The CodeFest has also been sponsored by \href{http://harbinger-partners.com/}{Harbinger Partners, Inc.}
and \href{https://curoverse.com/}{Curoverse} (the team behind the open source platform
\href{http://arvados.org/}{Arvados}).

\session{Keynote Speakers}{Days One \& Two, 9 \& 10 July 2014}{}

\section*{Keynote speaker: Philip Bourne (Day Two)}
\noindent
Dr. Bourne will speak about ``\textbf{Biomedical Research as an Open Digital Enterprise}'':
\begin{quote}
\small
\textit{The biomedical research lifecycle is fast becoming completely digital
and increasingly open to the point that publishing could simply become changing
the access control on given research objects comprising ideas, hypotheses, data,
software, results, conclusions, reviews, grants and so on. This offers immense
opportunities for software developers to enable the enterprise. I will describe
a vision for the digital enterprise and what the NIH and others are doing to
support the notion with the intent to accelerate scientific discovery.}
\end{quote}
Philip E. Bourne, PhD, is the Associate Director for Data Science at the NIH
and formerly Associate Vice Chancellor for Innovation and Industry Alliances,
and Professor of Pharmacology at the University of California San Diego.
Bourne's work at the NIH focuses on accelerating the rate of knowledge
discovery from the ever-increasing amounts of biomedical data at all scales --
from genomes to populations.
Bourne's laboratory focuses on relevant biological and educational outcomes
derived from computation and scholarly communication. This implies algorithms,
text mining, machine learning, metalanguages, biological databases, and
visualization applied to problems in systems pharmacology, evolution, cell
signaling, apoptosis, immunology and scientific dissemination. He has published
over 300 papers and 5 books. Bourne is the co-founder and founding Editor-in-Chief
of the open access journal PLOS Computational Biology. He is a Past President of
the International Society for Computational Biology, an elected fellow of the
American Association for the Advancement of Science (AAAS), the International
Society for Computational Biology (ISCB) and the American College of Medical
Informatics.
Bourne is committed to professional development through the Ten Simple Rules
series of articles and a variety of lectures and video presentations. Awards
include: the Jim Gray eScience Award (2010), the Benjamin Franklin Award (2009),
the Flinders University Convocation Medal for Outstanding Achievement (2004),
the Sun Microsystems Convergence Award (2002) and the CONNECT Award for new
inventions (1996 \& 97).

\section*{Keynote speaker: C.~Titus Brown (Day One)}
\noindent
Dr. Brown's topic is ``\textbf{A History of Bioinformatics (in the year 2039)}''.
\begin{quote}
\small
\textit{In 2039, I expect to look back at the last 25 years of biology and see
both wonderful surprises and missed opportunities. In this talk, I will attempt
to predict both some surprises and some of the opportunities I worry that we
will have missed over the next 25 years.}
\end{quote}
C.~Titus Brown is an assistant professor in the Department of Computer Science
and Engineering and the Department of Microbiology and Molecular Genetics. He
earned his PhD ('06) in developmental molecular biology from the California
Institute of Technology. Brown is director of the laboratory for Genomics,
Evolution, and Development (GED) at Michigan State University. He is a member
of the Python Software Foundation and an active contributor to the open source
software community. His research interests include computational biology,
bioinformatics, open source software development, and software engineering.

\newpage
\session{Genome-scale Data and Beyond}{Day One, 11 July 2014, 10:45 -- 12:30}{Session chair: Chris Fields}
\embed{2.5cm}{10cm}{2.5cm}{2.5cm}{GenomeScale-16-ADAM.pdf}
\embed{2.5cm}{10cm}{2.5cm}{2.5cm}{GenomeScale-32-Framework-RNASeq.pdf}
\embed{2.5cm}{5cm}{2cm}{2.5cm}{GenomeScale-27-SPADES.pdf}
\embed{2cm}{2.5cm}{1cm}{2.5cm}{GenomeScale-22-SigSeeker.pdf}
\embed{2.5cm}{3cm}{1.5cm}{2.5cm}{GenomeScale-21-Galaxy-extending.pdf}

\session{Visualization}{Day One, 11 July 2014, 14:00 -- 15:30}{Session chair: Rob Davey}

\embed{2.5cm}{2.75cm}{1cm}{2.5cm}{Visualization-39-WormGUIDES.pdf}
\embed{3cm}{10cm}{3cm}{2.5cm}{Visualization-05-BioJS.pdf}
\embed{2.5cm}{11cm}{2.5cm}{2cm}{Visualization-07-Biodalliance.pdf}
\embed{2cm}{2.5cm}{0cm}{2.5cm}{Visualization-17-TGAC-browser.pdf}
\embed{2.5cm}{10cm}{2.5cm}{2.5cm}{Visualization-38-IGB.pdf}

\session{Bioinformatics Open Source Project Updates}{Day One, 11 July 2014, 16:00 -- 17:00}{Session chair: Peter Cock}

\embed{2.5cm}{8cm}{2.5cm}{2.5cm}{Updates-20-BioMart.pdf}
\embed{2.5cm}{3cm}{1cm}{2.5cm}{Updates-25-Biocaml.pdf}
\embed{2cm}{2.5cm}{2cm}{3cm}{Updates-41-BioRuby.pdf}
\embed{2cm}{1.5cm}{0cm}{1cm}{Updates-19-Biopython.pdf}
\embed{2.25cm}{3cm}{1.5cm}{2.5cm}{Updates-02-UGENE.pdf}

\session{Software Interoperability}{Day Two, 12 July 2014, 10:45 -- 12:30}{Session chair: Raoul Bonnal}

\embed{2.5cm}{3cm}{1.5cm}{2.5cm}{Interoperability-04-Pathview.pdf}
\embed{2.5cm}{2.5cm}{1cm}{2.5cm}{Interoperability-13-Mobyle2.pdf}
\embed{2.5cm}{2.2cm}{1cm}{3cm}{Interoperability-34-OWL-computing.pdf}
\embed{2.5cm}{2.5cm}{0.5cm}{3.5cm}{Interoperability-14-Taverna-player.pdf}
\embed{2.5cm}{2.5cm}{1.5cm}{3cm}{Interoperability-40-Small-Tools.pdf}

\session{Open Science and Reproducible Research}{Day Two, 12 July 2014, 14:00 -- 15:30}{Session chair: Hilmar Lapp}

\embed{2.5cm}{10cm}{2cm}{2.5cm}{OpenScience-11-SEEK.pdf}
\embed{2.5cm}{5cm}{2.5cm}{2.5cm}{OpenScience-31-Arvados.pdf}
\embed{3cm}{2.5cm}{1cm}{2.5cm}{OpenScience-35-Galaxy-community.pdf}
\embed{2.5cm}{5cm}{2.5cm}{2.5cm}{OpenScience-23-Community-Pipelines.pdf}
\embed{2.5cm}{3.25cm}{2.5cm}{3cm}{OpenScience-26-Open-Tree-of-Life.pdf}

\end{document}
